\documentclass[11pt,a4paper]{article}
\usepackage[utf8]{inputenc}
\usepackage[T1]{fontenc}
\usepackage{lmodern}
\usepackage[margin=2.5cm]{geometry}
\usepackage{amsmath,amssymb,amsfonts}
\usepackage{graphicx}
\graphicspath{{FIGURAS_FINALES/}}
\usepackage{booktabs}
\usepackage{natbib}
\usepackage{hyperref}
\usepackage{float}
\usepackage{microtype}
\usepackage{tabularx}
\usepackage{xurl}
\usepackage{caption}
\usepackage{subcaption}

% Float placement parameters for optimal page layout and reduced whitespace
\renewcommand{\topfraction}{0.85}
\renewcommand{\bottomfraction}{0.85}
\renewcommand{\textfraction}{0.15}
\renewcommand{\floatpagefraction}{0.75}

\hypersetup{
    colorlinks=true,
    linkcolor=blue,
    citecolor=blue,
    urlcolor=blue
}

\title{\textbf{Multilinear and Subsurface Evaluation of the ENSO $E$ and $C$ Indices:\\ Quantifying Residual Dynamical Information, Temporal Irreversibility, and Threshold Sensitivity in the Equatorial Pacific}}

\author{M. Andrade\thanks{Corresponding author: \texttt{mandrade@igp.gob.pe}} \\
\small Direcci\'on de Ciencias de la Atm\'osfera , Hidr\'osfera y Cambio Clim\'atico, Instituto Geof\'isico del Per\'u (IGP),\\
Lima, Per\'u 
}

\date{November 2026}

\begin{document}

\maketitle

\begin{abstract}
The diversity of the El Ni\~no--Southern Oscillation (ENSO) is canonically characterized by the orthogonal Eastern ($E$) and Central ($C$) Pacific sea surface temperature indices. While empirical observations exhibit pronounced phase-space curvature, whether higher-order multilinear moment tensors provide genuinely novel physical information beyond linear wave kinematics, Takahashi's quadratic parameterization, and recharge-oscillator dynamics remains an open scientific question. Here, we present an exhaustive diagnostic assessment of the trivariate $(E, C, D20)$ system using monthly observations (1980--2018, $N = 466$ continuous months), incorporating eastern equatorial 20$^\circ$C isotherm depth ($D20$) and zonal wind stress. We show that the static third-order tensor is largely reconstructible from established statistical relationships, while its residual components do not, by themselves, establish independent physical mechanisms: linear wave projection of the thermocline onto $(E, C)$ accounts for $71.76\%$ of the squared Frobenius norm ($\|\mathcal{T}_{\mathrm{lin}}\|_F^2 / \|\mathcal{T}\|_F^2$), while exact Gram-Schmidt orthogonalization confirms that static mixed moments algebraically resolve into linear correlation and canonical parabolic geometry ($m_{\mathrm{res}} \equiv \gamma \mathrm{Var}(q) = 0.1367$). In contrast, the bi-lagged multilinear surface $T_{ECH}(\tau_1, \tau_2)$ provides statistical evidence of time-reversal asymmetry at physically motivated lags centered at $(\tau_1 = +3, \tau_2 = +2)$ months ($\Delta T = +0.2405$), reflecting precursor thermocline deepening and Central Pacific warming preceding eastern peak anomalies. While an unconstrained multi-lag exploratory search across 625 combinations cannot reject random field fluctuations without prior physical knowledge ($p = 0.436$), this asymmetry is statistically significant against Fourier phase-randomized surrogates ($p = 0.0105$) and dynamic recharge models ($p = 0.0140$) when constrained to equatorial wave-transit timescales, and persists across the modern 2000--2018 validation epoch. Jackknife resampling reveals high sensitivity to extreme episodes: excluding the three super El Ni\~no events ($9.66\%$ of the record) produces a $66.1\%$ reduction in the global Frobenius norm (and $58.3\%$ under re-standardization), reflecting strong leverage from extreme tails rather than continuous ambient multilinear turbulence, though historical sample size precludes proving an abrupt threshold bifurcation over continuous non-linearities. Finally, incorporating multilinear terms into predictive models yields no out-of-sample skill improvement, underscoring that higher-order moments serve as potent diagnostic markers of phase-space asymmetry rather than prognostic predictors.
\end{abstract}

\vspace{0.5em}
\noindent\textbf{Keywords:} ENSO diversity; E and C indices; Higher-order statistics; Multilinear moments; Thermocline dynamics; Recharge oscillator; Convective threshold; Surrogate data.

\vspace{1.0em}

\section{Introduction}

The El Ni\~no--Southern Oscillation (ENSO) is Earth's dominant mode of interannual climate variability, exerting widespread ecological and socio-economic impacts across the globe \citep{rasmusson1982variations, mcphaden1998tropical, cai2021changing}. Over the past two decades, extensive research has established that ENSO is not a single canonical oscillation, but exhibits pronounced spatial and temporal diversity \citep{capotondi2015understanding, timmermann2018enso}. In particular, warm events can mature predominantly in the Eastern Equatorial Pacific (EP El Ni\~no) or in the Central Pacific (CP El Ni\~no), each associated with distinct teleconnections, atmospheric convection patterns, and coastal impacts.

To objectively characterize this diversity, \citet{takahashi2011enso} introduced the orthogonal $E$ and $C$ indices, derived from the leading empirical orthogonal functions (EOF1 and EOF2) of tropical Pacific sea surface temperature (SST) anomalies. In this canonical framework, the $E$ index quantifies extreme warm anomalies concentrated along the South American coast and the far-eastern cold tongue, while the $C$ index reflects temperature anomalies centered near the dateline. Crucially, scatter plots in the $(E, C)$ phase plane reveal a pronounced non-linear geometry: positive $E$ anomalies (extreme EP events) are systematically accompanied by positive $C$ anomalies, forming a parabolic curve ($C \propto E^2$), whereas cold anomalies (La Ni\~na) are predominantly negative in $C$ with $E$ remaining close to zero \citep{takahashi2011enso, karamperidou2014intrinsic}.

Subsequent theoretical and modeling studies have illuminated the physical mechanisms responsible for this phase-space asymmetry. \citet{takahashi2016strong} and \citet{capotondi2020diversity} demonstrated that ENSO dynamics operate under distinct moderate and strong regimes governed by a non-linear atmospheric convective threshold: once eastern SSTs exceed approximately $28^\circ\text{C}$, deep atmospheric convection shifts eastward into the cold tongue, triggering a collapse of the equatorial trade winds and non-linear Kelvin wave amplification. Simultaneously, \citet{an2004nonlinearity} emphasized that oceanic non-linear dynamical heating (NDH)---specifically the horizontal and vertical advection of temperature anomalies by anomalous currents ($-\mathbf{u}' \cdot \nabla T'$)---acts as a powerful rectifying mechanism that intensifies warm extremes while damping cold anomalies.

Concurrently, higher-order statistical methods have been increasingly applied to elucidate non-Gaussianity in climate records \citep{monahan2004spatial, pires2021bispectral}. In bispectral theory, third-order moments and cumulants provide natural tools for detecting quadratic phase coupling and irreversible temporal evolution. In multilinear algebra, third-order moment tensors $\mathcal{T}^{(3)} \in \mathbb{R}^{d \times d \times d}$ have gained prominence for multi-way feature extraction and non-linear dimensionality reduction \citep{kolda2009tensor, comon2014tensors}.

However, a critical scientific question remains: \textit{Does the empirical formulation of a third-order moment tensor combining surface indices $(E, C)$ with subsurface thermocline depth ($D20$) provide genuinely novel physical information beyond known linear wave mechanics, Takahashi's quadratic parameterization, and documented threshold dynamics?}

The present work provides a rigorous evaluation of this question. Rather than asserting tensor novelty per se, we formulate the central contribution as: \textbf{identifying and quantifying the residual dynamical information of the trivariate $(E, C, D20)$ moments after removing the contributions from linear wave dependence, persistence, and known quadratic relationships.} Specifically, by analyzing synchronized monthly records of $E$, $C$, eastern equatorial thermocline depth ($D20$), and zonal wind stress ($\tau_x$) over 1980--2018 ($N = 466$ continuous months), we:
\begin{enumerate}
    \item Deconstruct the static third-order moment tensor $\mathcal{T}^{(3)}$ via linear thermocline wave projection and exact Gram-Schmidt orthogonalization;
    \item Quantify the temporal irreversibility of the bi-lagged multilinear moment surface $T_{ECH}(\tau_1, \tau_2)$ and validate its statistical significance against linear Gaussian surrogate data and coupled dynamic recharge models;
    \item Assess the leverage of extreme super El Ni\~no events via Jackknife subsampling to determine whether non-linearity represents continuous multilinear turbulence or episodic threshold activation; and
    \item Test out-of-sample predictive forecasts to evaluate whether empirical bilinear interaction terms provide measurable predictive skill beyond conventional quadratic models.
\end{enumerate}

\section{Data and Observational Timeline}

\subsection{Observational Datasets (1980--2018)}
Four monthly observational series were synchronized over the 39-year period March 1980 to December 2018 ($N = 466$ continuous months):
\begin{enumerate}
    \item \textbf{$E$ and $C$ Indices:} Derived from NOAA Extended Reconstructed Sea Surface Temperature version 5 (ERSSTv5; \citealt{huang2017extended}) over the tropical Pacific domain ($10^\circ\text{S}$--$10^\circ\text{N}$, $120^\circ\text{E}$--$80^\circ\text{W}$). Following the mathematical methodology of \citet{takahashi2011enso}, monthly SST anomalies relative to the 1981--2010 base period climatology were subjected to Empirical Orthogonal Function (EOF) analysis. The leading two principal component time series, $PC1(t)$ and $PC2(t)$, were rotated by $45^\circ$ according to the canonical definitions:
    \begin{equation}
    E(t) = \frac{PC1(t) - PC2(t)}{\sqrt{2}}, \qquad C(t) = \frac{PC1(t) + PC2(t)}{\sqrt{2}}
    \end{equation}
    High-frequency intraseasonal noise was removed using a standard 1-2-1 binomial filter. The spatial loading patterns and time series derived from ERSSTv5 correlate at $r > 0.98$ with the original HadISST-derived indices of \citet{takahashi2011enso}, confirming physical and mathematical equivalence.
    \item \textbf{Eastern Equatorial Thermocline Depth ($D20$):} Derived from the depth of the 20$^\circ$C isotherm averaged over the Eastern Pacific equatorial box ($2^\circ\text{S}$--$2^\circ\text{N}$, $110^\circ\text{W}$--$95^\circ\text{W}$), denoted $H(t)$, standardized relative to the 1981--2010 climatology. Data originate from the NOAA PMEL in situ equatorial moored buoy array compilation (\texttt{iso20\_complete.nc}) and the ECMWF Ocean Reanalysis System 5 (ORAS5; \citealt{zuo2019ecmwf}).
    \item \textbf{Atmospheric Zonal Wind Stress ($\tau_x$):} Extracted from ECMWF ORAS5 variable \texttt{sozotaux} over the Western-Central Pacific trigger region ($2^\circ\text{S}$--$2^\circ\text{N}$, $160^\circ\text{E}$--$160^\circ\text{W}$).
\end{enumerate}

\subsection{Subsurface Observational Provenance and Robustness Audit}
A rigorous assessment of subsurface data provenance across the observational timeline is essential:
\begin{itemize}
    \item \textbf{Pre-ATLAS Era (1980--1984):} Subsurface thermocline measurements in the eastern equatorial Pacific box ($110^\circ\text{W}$--$95^\circ\text{W}$) are grounded in in situ moorings from the NOAA Equatorial Pacific Ocean Climate Studies (EPOCS) program \citep{halpern1987observations} and historical CTD/XBT hydrographic transects. While \citet{halpern1987observations} documents the continuous deployment and moored temperature observations along the equator at $0^\circ, 110^\circ\text{W}$ and $0^\circ, 95^\circ\text{W}$ during 1980--1985, our direct data audit of the PMEL historical records indicates that in situ observations were available for $86.3\%$ of months at the $110^\circ\text{W}$ site during this 1980--1984 window. To ensure continuous temporal coverage and spatial coherence across data-sparse months prior to the full ATLAS deployment, these in situ observations were integrated and cross-validated with the continuous, dynamically constrained ECMWF ORAS5 ocean reanalysis, which assimilates all available historical subsurface temperature profiles using a 3D-Var assimilation scheme into the NEMO ocean model.
    \item \textbf{TOGA / ATLAS Expansion (1985--1994):} Following the 1982--83 El Ni\~no, autonomous ATLAS thermistor moorings were progressively deployed across the equatorial Pacific, providing continuous real-time telemetered thermal structure.
    \item \textbf{Mature Operational Array (1995--2018):} The full basin-wide 67-mooring TAO/TRITON array operated with $>95\%$ data completeness. Across the entire 466 months in the eastern box, only 4 monthly values ($0.86\%$ of the record) were missing in the in situ series (September--October 2012 and January--February 2013 during a TAO maintenance gap), which were linearly interpolated.
\end{itemize}

To ensure that results are not sensitive to pre-ATLAS hydrographic sampling, we conducted an epoch sensitivity test restricting all calculations to the modern ATLAS operational era (1985--2018, $N = 408$ continuous months). The linear regression coefficients between thermocline depth and surface indices are virtually identical: $b_E = 0.4811$ ($1985$--$2018$) vs. $0.4934$ ($1980$--$2018$); $b_C = 0.6223$ vs. $0.6156$; and $R^2 = 66.78\%$ vs. $66.54\%$. Third-order moments show identical stability: $T_{EEE} = 1.5632$ vs. $1.5708$ ($0.5\%$ change), $T_{EEH} = 0.9831$ vs. $1.0169$ ($3.3\%$ change), and $T_{ECH} = 0.3355$ vs. $0.3454$ ($2.9\%$ change), with all variations remaining strictly within $3.6\%$. This confirms that the statistical and physical conclusions are robust across observational epochs.

All four series were standardized by subtracting the sample mean and dividing by the sample standard deviation ($ddof = 1$):
\begin{equation}
e(t) = \frac{E(t) - \bar{E}}{\sigma_E}, \quad c(t) = \frac{C(t) - \bar{C}}{\sigma_C}, \quad h(t) = \frac{H(t) - \bar{H}}{\sigma_H}, \quad w(t) = \frac{\tau_x(t) - \bar{\tau}_x}{\sigma_{\tau_x}}
\end{equation}

\section{Mathematical and Methodological Framework}

\subsection{Trivariate Third-Order Moment Tensor $\mathcal{T}^{(3)}$}
Let $\mathbf{x}(t) = [x_1(t), x_2(t), x_3(t)]^T \equiv [e(t), c(t), h(t)]^T \in \mathbb{R}^3$. The static third-order moment tensor $\mathcal{T}^{(3)} \in \mathbb{R}^{3 \times 3 \times 3}$ is defined by:
\begin{equation}
\mathcal{T}_{ijk} = \frac{1}{N} \sum_{t=1}^N x_i(t) \, x_j(t) \, x_k(t), \quad i,j,k \in \{1, 2, 3\}
\end{equation}
Due to index permutation symmetry, $\mathcal{T}^{(3)}$ is fully supersymmetric and possesses exactly $\binom{3+3-1}{3} = 10$ independent components. Its magnitude is quantified by the Frobenius norm:
\begin{equation}
\|\mathcal{T}^{(3)}\|_F = \sqrt{\sum_{i=1}^3 \sum_{j=1}^3 \sum_{k=1}^3 \mathcal{T}_{ijk}^2}
\end{equation}

\subsection{Linear Wave Projection and Frobenius Norm Deconstruction}
Following equatorial shallow-water wave theory \citep{jin1997paradigm, meinen2000observations}, eastern thermocline depth responds as an integrated linear wave receiver forced by surface wind and thermal fields. We project $h(t)$ onto the standardized surface indices:
\begin{equation}
h(t) = a_1 e(t) + a_2 c(t) + \varepsilon_h(t)
\label{eq:ols_h}
\end{equation}
where $\mathbb{E}[\varepsilon_h e] = \mathbb{E}[\varepsilon_h c] = 0$. Substituting Eq.~(\ref{eq:ols_h}) into third-order moments involving $h$ yields:
\begin{equation}
\mathcal{T}_{ijh} = a_1 \mathcal{T}_{ije} + a_2 \mathcal{T}_{ijc} + \mathbb{E}[x_i x_j \varepsilon_h] \equiv \mathcal{T}_{ijh}^{\text{lin}} + \mathcal{T}_{ijh}^{\text{res}}
\end{equation}
The linear reconstructed tensor $\mathcal{T}_{\text{lin}}^{(3)}$ is defined by replacing $h$ with $h_{\text{lin}} = a_1 e + a_2 c$, leaving the true residual tensor:
\begin{equation}
\mathcal{T}_{\text{res}}^{(3)} = \mathcal{T}^{(3)} - \mathcal{T}_{\text{lin}}^{(3)}
\end{equation}
Importantly, because $\mathcal{T}_{\text{lin}}$ and $\mathcal{T}_{\text{res}}$ are not necessarily orthogonal in tensor Frobenius space, the squared Frobenius norm partitions exactly as:
\begin{equation}
\|\mathcal{T}^{(3)}\|_F^2 = \|\mathcal{T}_{\text{lin}}^{(3)} + \mathcal{T}_{\text{res}}^{(3)}\|_F^2 = \|\mathcal{T}_{\text{lin}}^{(3)}\|_F^2 + 2\langle \mathcal{T}_{\text{lin}}^{(3)}, \mathcal{T}_{\text{res}}^{(3)} \rangle_F + \|\mathcal{T}_{\text{res}}^{(3)}\|_F^2
\label{eq:frobenius_balance}
\end{equation}
where $\langle \mathcal{A}, \mathcal{B} \rangle_F = \sum_{ijk} \mathcal{A}_{ijk} \mathcal{B}_{ijk}$ is the Frobenius inner product. Therefore, while $\|\mathcal{T}_{\text{lin}}\|_F^2 / \|\mathcal{T}\|_F^2$ provides the squared Frobenius ratio, the cross-term $2\langle \mathcal{T}_{\text{lin}}, \mathcal{T}_{\text{res}} \rangle_F$ must be explicitly reported to document the geometric non-orthogonality of the multilinear decomposition.

\subsection{Gram-Schmidt Algebraic Partitioning of the Parabolic Term}
To evaluate the mixed surface moment $T_{EEC} = \mathbb{E}[e^2 c]$, we apply the Gram-Schmidt orthogonalization procedure to decouple $e(t)^2$ from the linear basis $\{1, e(t)\}$. Defining the general orthogonalized quadratic variable:
\begin{equation}
q(t) = e(t)^2 - \mathbb{E}[e^2] - \frac{\mathbb{E}[e^3]}{\mathbb{E}[e^2]} \, e(t)
\label{eq:gs_q}
\end{equation}
which identically satisfies $\mathbb{E}[q] = 0$ and $\text{Cov}(q, e) = \mathbb{E}[q e] = 0$ for any sample normalization convention ($ddof = 0$ or $ddof = 1$). Its exact variance is:
\begin{equation}
\text{Var}(q) = \mathbb{E}[e^4] - (\mathbb{E}[e^2])^2 - \frac{(\mathbb{E}[e^3])^2}{\mathbb{E}[e^2]}
\end{equation}
Regressing $c(t)$ onto the orthogonal basis $\{e(t), q(t)\}$:
\begin{equation}
c(t) = r_{\text{lin}} e(t) + \gamma q(t) + \eta(t)
\end{equation}
where $r_{\text{lin}} = \mathbb{E}[e c] / \mathbb{E}[e^2]$ and $\gamma = \text{Cov}(c, q) / \text{Var}(q)$. For the linear residual $c_{\text{res}}(t) = c(t) - r_{\text{lin}} e(t)$, the residual cross-moment satisfies the exact algebraic identity:
\begin{equation}
m_{\text{res}} \equiv \mathbb{E}[e^2 c_{\text{res}}] = T_{EEC} - r_{\text{lin}} \mathbb{E}[e^3] \equiv \gamma \cdot \text{Var}(q)
\label{eq:gs_identity}
\end{equation}
\textbf{Physical and Algebraic Meaning:} Eq.~(\ref{eq:gs_identity}) is an exact algebraic consequence of projecting onto the polynomial basis:
\begin{equation}
T_{EEC} = \underbrace{r_{\text{lin}} \, T_{EEE}}_{\text{Linear Projection}} + \underbrace{\gamma \cdot \text{Var}(q)}_{\text{Parabolic Curvature}}
\end{equation}
This confirms that in static surface data, the third-order cross-moment $T_{EEC}$ contains no residual component orthogonal to Takahashi's parabolic regression. Note, however, that this identity represents an exact algebraic decomposition guaranteed by mathematical construction; it does not rule out cubic or higher-order non-linear dynamics in the time-dependent prognostic equations of the climate system (such as nonlinear cubic damping in the equatorial ocean heat recharge; cf. \citealt{han2026realistic}).

\subsection{Bi-Lagged Cross-Moments, Temporal Irreversibility, and Surrogate Testing}
To investigate dynamical time-lagged relationships, we define the bi-lagged third-order cross-moment surface:
\begin{equation}
T_{ECH}(\tau_1, \tau_2) = \frac{1}{N_{\tau}} \sum_t e(t) \, c(t - \tau_1) \, h(t - \tau_2)
\label{eq:bilagged}
\end{equation}
over discrete lags $\tau_1, \tau_2 \in [-12, +12]$ months, where $N_{\tau}$ is the number of valid overlapping months. Under this lag convention:
\begin{itemize}
    \item $\tau_1 > 0$ indicates that the Central Pacific anomaly $c$ occurs $\tau_1$ months \textbf{before} eastern warming $e(t)$ ($C$ leads $E$).
    \item $\tau_2 > 0$ indicates that eastern thermocline deepening $h$ occurs $\tau_2$ months \textbf{before} eastern warming $e(t)$ ($D20$ leads $E$).
\end{itemize}
To rigorously assess temporal irreversibility (the dynamical arrow of time), we evaluate the time-reversal asymmetry metric around the identified dynamical peak $(\tau_1^*, \tau_2^*)$ \citep{pires2021bispectral}:
\begin{equation}
\Delta T_{ECH}(\tau_1^*, \tau_2^*) = T_{ECH}(\tau_1^*, \tau_2^*) - T_{ECH}(-\tau_1^*, -\tau_2^*)
\label{eq:asymmetry_metric}
\end{equation}
To determine whether $\Delta T_{ECH}$ reflects true non-linear dynamics or is a sampling artifact of linear persistence or known static curvature, we establish four statistical benchmarks:
\begin{enumerate}
    \item \textbf{Moving Block Bootstrap:} Resampling blocks of length $L = 18$ months ($B = 2000$ iterations) to construct non-parametric 95\% confidence intervals preserving local persistence.
    \item \textbf{Multivariate Fourier Phase-Randomized Surrogates:} Generating $B = 2000$ surrogate realizations by applying identical random phase rotations to the Fourier transform across channels, perfectly preserving power spectra, autocorrelations, and cross-covariances while destroying higher-order non-linear coupling.
    \item \textbf{Trivariate Gaussian VAR(2) Null Model:} Fitting a stationary vector autoregression of order 2 to $[e, c, h]^T$ and simulating $B = 2000$ synthetic realizations of length $N = 466$.
    \item \textbf{Coupled Non-Linear Quadratic Recharge Null Model:} Simulating a bivariate VAR(2) for $[E, H]$ coupled with Takahashi's dynamic quadratic relation $C(t) = c_0 + \rho_1 C(t-1) + \rho_2 C(t-2) + \beta_0 E(t) + \beta_1 E(t-1) + \gamma E(t)^2 + \theta_0 H(t) + \theta_1 H(t-1) + \varepsilon_C(t)$ ($B = 1000$). This benchmark explicitly preserves linear wave dynamics, thermal feedback, and the high temporal autocorrelation of $C$ ($r = 0.978$). We evaluate both the a priori hypothesis test at the dynamically predicted lag $(\tau_1 = +3, \tau_2 = +2)$ and the global max-statistic across all 625 combinations in the $[-12, +12]^2$ grid to assess the look-elsewhere effect.
\end{enumerate}

\subsection{Out-of-Sample Predictive Framework}
We formulate five nested predictive models for forecasting $E(t+k)$ at lead times $k \in \{1, 3, 6\}$ months:
\begin{itemize}
    \item \textbf{M1 [Autoregressive AR(1)]:} $E(t+k) = \beta_0 + \beta_1 E(t)$.
    \item \textbf{M2 [Bivariate Linear]:} $E(t+k) = \beta_0 + \beta_1 E(t) + \beta_2 C(t)$.
    \item \textbf{M3 [Physical Linear]:} $E(t+k) = \beta_0 + \beta_1 E(t) + \beta_2 C(t) + \beta_3 H(t) + \beta_4 \tau_x(t)$.
    \item \textbf{M4 [Quadratic]:} $M3 + \beta_5 E(t)^2 + \beta_6 C(t)^2$.
    \item \textbf{M5 [Bilinear Interaction]:} $M4 + \beta_7 E(t)H(t) + \beta_8 C(t)H(t) + \beta_9 H(t)^2$.
\end{itemize}
Models are trained on 1980--2005 ($N_{\text{train}} = 310$ months) and tested out-of-sample on 2006--2018 ($N_{\text{test}} = 156$ months), with standardization statistics computed strictly on the training partition. Forecast accuracy differences are evaluated using RMSE and tested via the \citet{diebold1995comparing} test with Harvey-Leybourne-Newbold small-sample correction and Bartlett serial error weighting.

\section{Results}

\subsection{Observational Covariation and Phase-Space Geometry}

\begin{figure}[!htbp]
\centering
\includegraphics[width=\textwidth]{fig01_joint_distributions_ECD20.pdf}
\caption{Joint scatter distributions, quadratic and linear regression fits, and marginal univariate histograms of standardized indices $e(t)$, $c(t)$, and $h(t)$ ($N = 466$, 1980--2018). (a) $(e, c)$ plane colored by thermocline depth anomaly $h$, showing the quadratic regression curve $C(E)$ with 95\% confidence interval for the fitted conditional mean. (b) $e$ vs. $h$ colored by $c$ with linear fit ($r = +0.537$). (c) $c$ vs. $h$ colored by $e$ with linear fit ($r = +0.651$). (d) Compared marginal probability densities with empirical skewness values.}
\label{fig:joint_dist}
\end{figure}

Figure~\ref{fig:joint_dist} displays the bivariate distributions among $e$, $c$, and $h$. The $(e, c)$ plane exhibits the characteristic parabolic asymmetry documented by \citet{takahashi2011enso}. In contrast, the relationships involving the thermocline ($e, h$ and $c, h$) are strongly linear. The bivariate OLS regression of eastern thermocline depth onto surface indices yields:
\begin{equation}
h(t) = 0.4934 \, e(t) + 0.6156 \, c(t) + \varepsilon_h(t), \quad R^2 = 66.54\%
\label{eq:ols_empirical}
\end{equation}
confirming that approximately two-thirds of thermocline depth variance is linearly explained by surface temperatures.

\subsection{Static Tensor Deconstruction and Gram-Schmidt Proof}

Table~\ref{tab:tensor} presents the 10 independent components of $\mathcal{T}^{(3)}$, contrasted against the linear statistical reconstruction based on Eq.~(\ref{eq:ols_empirical}).

\begin{table}[!htbp]
\centering
\small
\caption{Independent components of the third-order moment tensor $\mathcal{T}^{(3)}$ ($N=466$, 1980--2018; standardized with $ddof=1$). All values and multiplicities exactly reproduce the total Frobenius norm and projection identities.}
\label{tab:tensor}
\resizebox{\linewidth}{!}{%
\begin{tabular}{llcccc}
\toprule
Component & Notation & Observed & Linear Reconstruction & Residual & Linear Expl. \% \\
\midrule
$\mathcal{T}_{111}$ & $T_{EEE}$ & \textbf{+1.5708} & --- & --- & Univariate Skewness \\
$\mathcal{T}_{222}$ & $T_{CCC}$ & \textbf{-0.4430} & --- & --- & Univariate Skewness \\
$\mathcal{T}_{333}$ & $T_{HHH}$ & \textbf{+0.8682} & +0.3295 & +0.5386 & 38.0\% \\
$\mathcal{T}_{112}$ & $T_{EEC}$ & \textbf{+0.2478} & +0.2478 & +0.0000 & 100.0\% (GS Identity) \\
$\mathcal{T}_{122}$ & $T_{ECC}$ & \textbf{+0.2369} & +0.2369 & +0.0000 & 100.0\% (Surface) \\
$\mathcal{T}_{113}$ & $T_{EEH}$ & \textbf{+1.0169} & \textbf{+0.9275} & \textbf{+0.0894} & \textbf{91.2\%} \\
$\mathcal{T}_{223}$ & $T_{CCH}$ & \textbf{-0.1649} & \textbf{-0.1559} & \textbf{-0.0090} & \textbf{94.5\%} \\
$\mathcal{T}_{133}$ & $T_{EHH}$ & \textbf{+0.9188} & +0.6227 & +0.2962 & 67.8\% \\
$\mathcal{T}_{233}$ & $T_{CHH}$ & \textbf{+0.1603} & +0.0363 & +0.1240 & 22.7\% \\
$\mathcal{T}_{123}$ & $T_{ECH}$ & \textbf{+0.3454} & \textbf{+0.2681} & \textbf{+0.0773} & \textbf{77.6\%} \\
\bottomrule
\end{tabular}%
}
\end{table}

The total Frobenius norm of the empirical tensor is $\|\mathcal{T}^{(3)}\|_F = 3.2061$ ($\|\mathcal{T}^{(3)}\|_F^2 = 10.2793$)\footnote{Under population standardization ($ddof=0$), $\|\mathcal{T}\|_F = 3.2165$, $\|\mathcal{T}_{\text{lin}}\|_F = 2.7248$, $\|\mathcal{T}_{\text{res}}\|_F = 0.8147$, $T_{EEE} = 1.5758$, $T_{CCC} = -0.4444$, $T_{EEH} = 1.0202$, and $T_{ECH} = 0.3465$, which differ strictly by the constant factor $(N/(N-1))^{3/2} = 1.003228$ ($0.32\%$). The ratio of squared Frobenius norms is invariant ($71.76\%$). Both the 3D tensor array and the weighted sum over the 10 components with their exact multinomial multiplicities yield identical norms to machine precision ($0.00 \times 10^0$).}. The linear statistical reconstruction tensor possesses a norm of $\|\mathcal{T}_{\text{lin}}^{(3)}\|_F = 2.7160$ ($\|\mathcal{T}_{\text{lin}}^{(3)}\|_F^2 = 7.3768$), which accounts for \textbf{71.76\%} of the total squared Frobenius norm ($\|\mathcal{T}_{\text{lin}}\|_F^2 / \|\mathcal{T}\|_F^2$).

Because the reconstructed tensor $\mathcal{T}_{\text{lin}}$ and the residual tensor $\mathcal{T}_{\text{res}}$ are not orthogonal in tensor space, evaluating the exact quadratic balance (Eq.~\ref{eq:frobenius_balance}) yields:
\begin{equation}
\underbrace{10.2793}_{\|\mathcal{T}\|_F^2 \, (100\%)} = \underbrace{7.3768}_{\|\mathcal{T}_{\text{lin}}\|_F^2 \, (71.76\%)} + \underbrace{2.2431}_{2\langle\mathcal{T}_{\text{lin}}, \mathcal{T}_{\text{res}}\rangle_F \, (21.82\%)} + \underbrace{0.6594}_{\|\mathcal{T}_{\text{res}}\|_F^2 \, (6.42\%)}
\end{equation}
Importantly, the percentages $71.76\%$, $21.82\%$, and $6.42\%$ represent algebraic terms of this exact norm balance rather than three orthogonal fractions of explained variance. The true residual tensor norm is $\|\mathcal{T}_{\text{res}}^{(3)}\|_F = \sqrt{0.6594} = 0.8121$, representing a residual ratio of $25.33\%$ ($\|\mathcal{T}_{\text{res}}\|_F / \|\mathcal{T}\|_F$). Crucially, for mixed moments involving the thermocline, linear statistical coupling accounts for $91.2\%$ of $T_{EEH}$, $94.5\%$ of $T_{CCH}$, and $77.6\%$ of $T_{ECH}$. Evaluating the theoretical projection identity $T_{ECH}^{\text{lin}} = a_1 T_{EEC} + a_2 T_{ECC} = 0.4934(0.2478) + 0.6156(0.2369) = \mathbf{+0.2681}$ confirms exact consistency between the regression coefficients and the empirical moments to machine precision ($1.11 \times 10^{-16}$).

\begin{figure}[!htbp]
\centering
\includegraphics[width=\textwidth]{fig02_linear_vs_nonlinear_decomposition.pdf}
\caption{Linear statistical reconstruction and residual third-order moment partitioning across the components of $\mathcal{T}^{(3)}$ (matching Table~\ref{tab:tensor}). (a) Moment partition showing observed moments $T_{ijk}$, linear statistical reconstructions $T_{ijk}^{\text{lin}}$, and residual third-order moments $T_{ijk}^{\text{res}}$, with linear reconstruction percentages indicated for key thermocline cross-moments ($T_{EEH}: 91.2\%$, $T_{CCH}: 94.5\%$, $T_{ECH}: 77.6\%$). (b) Frobenius squared norm balance ($\|\mathcal{T}\|_F^2 = 10.2793$, where $71.76\%$ denotes the ratio of squared Frobenius norms $\|\mathcal{T}_{\text{lin}}\|_F^2 / \|\mathcal{T}\|_F^2$, $21.82\%$ the non-orthogonal cross-term, and $6.42\%$ the residual norm ratio $\|\mathcal{T}_{\text{res}}\|_F^2 / \|\mathcal{T}\|_F^2$) and verification of the Gram-Schmidt algebraic orthogonality identity on $T_{EEC}$ ($m_{\text{res}} = \gamma \text{Var}(q) = 0.1367$, $\Delta = 0.00$).}
\label{fig:decomp}
\end{figure}

Evaluating the general Gram-Schmidt identity (Eq.~\ref{eq:gs_identity}) for the bivariate moment $T_{EEC}$:
\begin{align*}
\mathbb{E}[e^2] &= 0.9979, \quad \text{Var}(q) = 2.5554, \quad \gamma = 0.0535 \implies \gamma \cdot \text{Var}(q) = 0.1367 \\
m_{\text{res}} &= T_{EEC} - r_{\text{lin}} T_{EEE} = 0.2478 - (0.0708)(1.5708) = 0.1367
\end{align*}
The difference $|m_{\text{res}} - \gamma \text{Var}(q)| = 0.00 \times 10^{00}$ to machine precision. As established in Section~3.3, this represents an exact algebraic decomposition of the selected third-order cross-moment, guaranteeing orthogonality to the polynomial basis by mathematical construction rather than proving that nature lacks alternative non-linear dynamics.

\subsection{Temporal Irreversibility and the Arrow of Time}

\begin{figure}[!htbp]
\centering
\includegraphics[width=\textwidth]{fig03_lagged_tensor_surface_ECH.pdf}
\caption{Bi-lagged cross-third moment surface and statistical evaluation of temporal irreversibility. (a) Bi-lagged surface $T_{ECH}(\tau_1, \tau_2) = \mathbb{E}[E(t) \, C(t - \tau_1) \, D20(t - \tau_2)]$, displaying the global maximum at $(\tau_1 = +3, \tau_2 = +2)$ months ($+0.4293$) and its time-reversed counterpart at $(-3, -2)$ ($+0.1888$). (b) Time-reversal asymmetry map $\Delta T_{ECH}(\tau_1, \tau_2) = T - T_{\text{rev}}$, highlighting the asymmetry evaluated at the selected $T$-peak ($\Delta T(+3, +2) = +0.2405$, yellow circle), the unconstrained global maximum of asymmetry at $(\tau_1 = +10, \tau_2 = +1)$ ($\Delta T = +0.6812$, magenta diamond), and the physical wave-transit window $\Omega_{\text{phys}}$ (dashed rectangle, positive everywhere, mean $+0.2318$). (c) Monte Carlo null distributions for the targeted physical lag (Fourier surrogates, $p=0.0105$; Dynamic Coupled Quadratic, $p=0.0140$) contrasted against the unconstrained 625-lag exploratory search distribution ($p=0.4360$, testing asymmetry at the argmax of $T_{ECH}$ without physical priors). (d) Retrospective temporal validation on the modern era (2000--2018, $N=228$), showing the moving block bootstrap distribution ($L=18$, $B=2000$; 95\% CI $[+0.0319, +0.5125]$) and rejection of the 2000--2018 dynamic quadratic null model ($p=0.0065$ unilateral, $p=0.0150$ bilateral).}
\label{fig:lagged}
\end{figure}

Figure~\ref{fig:lagged} displays the bi-lagged cross-moment surface $T_{ECH}(\tau_1, \tau_2)$. In contrast to the static domain, a distinct dynamical signature emerges:
\begin{itemize}
    \item A global maximum of the third-order cross-moment \textbf{$+0.4293$} occurs at:
    \begin{align}
    \tau_1^* &= +3 \text{ months (Central Pacific precursor leads } E(t)), \nonumber \\
    \tau_2^* &= +2 \text{ months (Thermocline depth } D20 \text{ leads } E(t))
    \end{align}
    This alignment captures the canonical sequence of El Ni\~no development: Central Pacific warming ($\tau_1 = +3$) and westerly wind stress force eastward downwelling Kelvin waves, deepening the eastern thermocline ($\tau_2 = +2$), which subsequently triggers peak eastern SST warming ($E$, month 0).
    \item The time-reversed counterpart under $(t \to -t)$ is evaluated at $(-\tau_1^*, -\tau_2^*) = (-3, -2)$:
    \begin{equation}
    T_{ECH}(-3, -2) = \mathbb{E}[E(t) \, C(t+3) \, D20(t+2)] = +0.1888
    \end{equation}
    \item Evaluating the time-reversal asymmetry metric (Eq.~\ref{eq:asymmetry_metric}) at this identified $T$-peak:
    \begin{equation}
    \Delta T_{ECH}(+3, +2) = T_{ECH}(+3, +2) - T_{ECH}(-3, -2) = 0.4293 - 0.1888 = \mathbf{+0.2405}
    \end{equation}
    It should be noted that $+0.2405$ is the asymmetry evaluated specifically at the physical maximum of $T_{ECH}$. Across the entire unconstrained $25 \times 25 = 625$ lag grid, the mathematical maximum of $\Delta T$ occurs at $(\tau_1 = +10, \tau_2 = +1)$ ($\Delta T = +0.6812$), driven by the negative reversed tail ($T = +0.2366, T_{\text{rev}} = -0.4446$, reflecting multi-season phase reversal into La Ni\~na). This crucial distinction reinforces the necessity of our dual hypothesis testing framework.
\end{itemize}

To establish whether this asymmetry is statistically significant and represents novel dynamical information beyond known physical processes, we contrast $\Delta T_{ECH}(+3, +2)$ against four benchmark tests:
\begin{enumerate}
    \item \textbf{Moving Block Bootstrap ($B = 2000, L = 18$ months):} Yields a bootstrap mean of $+0.2081 \pm 0.0887$, with a 95\% confidence interval of $[+0.0430, +0.3872]$, confirming that the empirical asymmetry is strictly positive and bounded away from zero.
    \item \textbf{Multivariate Fourier Phase-Randomized Surrogates ($B = 2000$):} Surrogates preserving the exact power spectra, autocorrelations, and cross-covariances produce a null distribution of $-0.0024 \pm 0.0937$. The observed asymmetry of $+0.2405$ rejects linear Gaussian stationarity with \textbf{$p = 0.0105$} ($p < 0.05$).
    \item \textbf{Trivariate Gaussian VAR(2) Null Model ($B = 2000$):} Simulating a stationary linear Markovian process yields a null distribution of $+0.0012 \pm 0.1083$, rejecting the linear VAR benchmark with \textbf{$p = 0.0295$} ($p < 0.05$).
    \item \textbf{Dynamic Non-Linear Coupled Recharge Null Model ($B = 1000$):} Preserving the high persistence of $C$ ($r = 0.978$) and dynamic coupling yields a null distribution of $+0.0444 \pm 0.0941$ (95th percentile $= 0.1954$). The observed asymmetry of $+0.2405$ rejects this dynamic quadratic benchmark with \textbf{$p = 0.0140$} ($p < 0.05$).
\end{enumerate}

\textbf{Dual Hypothesis Framework: Multiple Testing vs. Physical Priors:}
To evaluate the impact of multiple testing across the $25 \times 25 = 625$ lag combinations in the $[-12, +12]^2$ grid, we conducted two complementary experiments:
\begin{enumerate}
    \item \textit{Experiment A (Unconstrained Multi-Lag Exploratory Search):} When the peak lag is selected dynamically inside each null simulation without prior physical constraints (evaluating asymmetry at the dynamically identified argmax of $T_{ECH}$ across all 625 combinations), the null distribution has a mean of $+0.2363 \pm 0.1221$ (95th percentile $= 0.4549$). Testing $\Delta T_{\text{obs}} = +0.2405$ against this unconstrained distribution yields \textbf{$p = 0.4360$} ($\approx 0.44$). This openly acknowledges that an unconstrained exploratory search across 625 lags cannot distinguish the observed asymmetry at the argmax of $T_{ECH}$ from extreme-value sampling noise of a persistent random field without physical wave-transit priors.
    \item \textit{Experiment B (Retrospective Temporal Cross-Validation and Physical Window):} In contrast, evaluating split-sample cross-validation between the calibration era (1980--1999, $N = 238$) and the subsequent modern era (2000--2018, $N = 228$) reveals that the physical lag $(\tau_1 = +3, \tau_2 = +2)$ remains robustly positive: in 2000--2018, $\Delta T_{ECH}(+3, +2) = \mathbf{+0.3227}$ ($T_{\text{fwd}} = 0.4908, T_{\text{rev}} = 0.1681$). Moving block bootstrap specifically over 2000--2018 ($L=18$ mo, $B=2000$) yields a 95\% confidence interval of $[+0.0319, +0.5125]$ (mean $+0.2602 \pm 0.1245$), strictly bounded away from zero. Sensitivity testing across block lengths $L \in \{12, 18, 24\}$ months confirms this stability (e.g., $95\%$ CI of $[+0.0419, +0.5248]$ for $L=24$ mo in 2000--2018, and $[+0.0807, +0.3543]$ for $L=24$ mo over the full record). Phase-randomized Fourier surrogates calculated specifically on 2000--2018 reject linear Gaussian dynamics at \textbf{$p = 0.0395$} ($p < 0.05$). Crucially, to test whether this persistence can be explained by known nonlinearities, we calibrated the dynamic coupled quadratic null model specifically over the 2000--2018 period ($B=2000$). The resulting null distribution has a mean of $+0.0029 \pm 0.1301$ (95th percentile $= 0.2142$), rejecting the dynamic quadratic benchmark with \textbf{$p = 0.0065$} (one-tailed, $p < 0.01$) and \textbf{$p = 0.0150$} (two-tailed, $p < 0.05$). While the full record was historically inspected, this retrospective temporal validation confirms that the surface-thermocline asymmetry is not an artifact of pre-2000 data or extreme 1982/1997 events, and provides evidence of temporal irreversibility beyond stationary linear or coupled quadratic models. In addition, evaluating the physical wave-transit window $\Omega_{\text{phys}} = \{\tau_1 \in [2,4], \tau_2 \in [1,3]\}$ yields strictly positive asymmetry across all 9 grid points (mean $= +0.2318$). Constrained to this physical wave hypothesis, the asymmetry is statistically established ($p = 0.0140$).
\end{enumerate}

\subsection{Physical Consistency with Oceanic Kelvin Waves and Wind Stress}

\begin{figure}[!htbp]
\centering
\includegraphics[width=\textwidth]{fig04_hovmoller_wind_thermocline_coupling.pdf}
\caption{(a) Equatorial Hovm\"oller diagram ($135^\circ$E--$90^\circ$W) of TAO/TRITON $D20$ anomalies during the 1997--98 El Ni\~no, with Kelvin wave propagation trajectory line ($c \approx 2.5\text{ m/s}$). (b) Equatorial zonal wind stress anomalies ($\tau_x$) from ORAS5. (c) Third-order mixed cross-moment $\mathbb{E}[E(t)H(t) \cdot \tau_x(t-\tau)]$, showing peak at $\tau = +1$ month. (d) Linear cross-correlation functions confirming that wind stress leads $D20$ by $+1$ month ($r=0.69$) and eastern SST ($E$) by $+2$ months ($r=0.40$).}
\label{fig:hovmoller}
\end{figure}

Figure~\ref{fig:hovmoller} confirms the physical wave mechanisms governing the subsurface. Western-central Pacific zonal wind stress anomalies ($\tau_x$) lead eastern thermocline depth ($D20$) by \textbf{$+1$ month} ($r = 0.69$), consistent with equatorial Kelvin wave transit timescales \citep{meinen2000observations}. In turn, thermocline deepening leads eastern surface warming ($E$) by \textbf{$+1$ to $+2$ months} ($r = 0.56$). Thus, subsurface thermocline depth acts primarily as a linear wave integrator, while non-linearities are localized at the air-sea boundary and in mixed-layer thermal processes.

\subsection{Predictive Out-of-Sample Falsification}

Table~\ref{tab:models} and Figure~\ref{fig:models} display out-of-sample forecast evaluations across lead times $k = 1, 3, 6$ months for the test period 2006--2018 ($N_{\text{test}} = 156$; trained on 1980--2005, $N_{\text{train}} = 310$).

\begin{table}[!htbp]
\centering
\small
\caption{Out-of-sample root mean square error (RMSE) for predicting $E(t+k)$ (test: 2006--2018, $N=156$; train: 1980--2005, $N=310$). Differences between M4 and M5 are statistically insignificant across all horizons under the Diebold-Mariano test.}
\label{tab:models}
\resizebox{\linewidth}{!}{%
\begin{tabular}{llccc}
\toprule
Model Class & Predictor Set & Lead $k=1$ mo & Lead $k=3$ mo & Lead $k=6$ mo \\
\midrule
M1: AR(1) Persistence & $E$ & 0.2447 & 0.6559 & 0.9422 \\
M2: Bivariate Linear & $E, C$ & 0.2285 & 0.5898 & 0.8149 \\
M3: Physical Linear & $E, C, H, \tau_x$ & 0.2112 & 0.5284 & 0.7876 \\
M4: Quadratic & $M3 + E^2, C^2$ & \textbf{0.2098} & \textbf{0.5215} & \textbf{0.7816} \\
M5: Bilinear & $M4 + E{\cdot}H, C{\cdot}H, H^2$ & 0.2115 (+0.81\%) & 0.5299 (+1.61\%) & 0.7958 (+1.82\%) \\
\bottomrule
\end{tabular}%
}
\end{table}

\begin{figure}[!htbp]
\centering
\includegraphics[width=\textwidth]{fig05_predictive_evaluation_out_of_sample.pdf}
\caption{Out-of-sample predictive evaluation across lead horizons $k = 1, 3, 6$ months (test: 2006--2018; train: 1980--2005). (a) Predictive RMSE for all five models from Table~\ref{tab:models}: persistence M1, bivariate linear M2, physical linear M3, conventional quadratic M4, and bilinear tensor M5. (b) Relative RMSE percentage change of bilinear model M5 over quadratic model M4 (+0.81\%, +1.61\%, +1.82\%). (c) Diebold-Mariano test statistics and $p$-values confirming no statistically significant difference from zero at $\alpha = 0.05$. (d) Squared error loss difference $(e_{\text{M5}}^2 - e_{\text{M4}}^2)$ over the verification period at lead $k=6$, displaying positive loss differences during 2015--16 without statistically significant forecast gain elsewhere.}
\label{fig:models}
\end{figure}

Incorporating physical predictors ($H, \tau_x$) and quadratic terms (M4) significantly improves forecast accuracy over persistence, reducing RMSE from 0.6559 to 0.5215 at lead 3. However, adding bilinear interaction terms ($E \cdot H, C \cdot H, H^2$; model M5) yields no incremental skill:
\begin{itemize}
    \item At lead 1 month, RMSE increases from 0.2098 to 0.2115 (\textbf{+0.81\%}).
    \item At lead 3 months, RMSE increases from 0.5215 to 0.5299 (\textbf{+1.61\%}).
    \item At lead 6 months, RMSE increases from 0.7816 to 0.7958 (\textbf{+1.82\%}).
\end{itemize}
In predictive regression, the interaction term $E \cdot H$ is not statistically significant ($p = 0.22 > 0.05$), and the \citet{diebold1995comparing} test confirms that forecast differences between M4 and M5 are statistically indistinguishable ($p = 0.66$ at $k=1$, $p = 0.66$ at $k=3$, and $p = 0.07$ at $k=6$). This indicates that adding empirical bilinear interaction terms yields no measurable forecast benefit in out-of-sample testing, producing slightly higher RMSE across all horizons due to parameter penalty without physical predictive gain.

\subsection{Jackknife Sensitivity: Extreme Event Concentration}

\begin{table}[!htbp]
\centering
\small
\caption{Jackknife sensitivity analysis of $\mathcal{T}^{(3)}$ to super El Ni\~no events ($N=466$, 1980--2018). Case 1 retains fixed global standardization ($\sigma$ from $N=466$) to measure extreme tail contributions directly. In Case 2 (subsamples re-standardized to unit variance), excluding all three super events yields $T_{EEE} = +0.4159$, $T_{EEH} = +0.1063$, $T_{ECH} = +0.2383$, and $\|\mathcal{T}^{(3)}\|_F = 1.3356$ ($-58.3\%$).}
\label{tab:jackknife}
\resizebox{\linewidth}{!}{%
\begin{tabular}{lccccc}
\toprule
Sample Selection & $T_{EEE}$ & $T_{EEH}$ & $T_{ECH}$ & $\|\mathcal{T}^{(3)}\|_F$ & Norm Change \\
\midrule
\textbf{Full Record (1980--2018)} & \textbf{+1.5708} & \textbf{+1.0169} & \textbf{+0.3454} & \textbf{3.2061} & Baseline \\
Excl. 1982--83 (16 mos; $N=450$) & +1.0223 & +0.6761 & +0.2675 & 2.1536 & -32.8\% \\
Excl. 1997--98 (15 mos; $N=451$) & +0.5532 & +0.4068 & +0.2209 & 1.5776 & -50.8\% \\
Excl. 2015--16 (14 mos; $N=452$) & +1.5014 & +0.9211 & +0.2262 & 2.8976 & -9.6\% \\
\textbf{Excl. All 3 Super Events (45 mos; $N=421$)} & \textbf{-0.1800} & \textbf{-0.1038} & \textbf{+0.0009} & \textbf{1.0873} & \textbf{-66.1\%} \\
\bottomrule
\end{tabular}%
}
\end{table}

\begin{figure}[!htbp]
\centering
\includegraphics[width=\textwidth]{fig06_jackknife_extreme_event_collapse.pdf}
\caption{Jackknife sensitivity and extreme event concentration in third-order statistics. (a) Moment stability of $T_{EEE}$, $T_{EEH}$, and $T_{ECH}$ under individual and combined exclusion of the three super El Ni\~no events. (b) Frobenius norm collapse comparison between global standardization (Case 1: $-66.1\%$) and subsample re-standardization (Case 2: $-58.3\%$). (c) Cumulative contribution to the third-order moment $T_{EEH} = \frac{1}{N}\sum_t E(t)^2 D20(t)$ over 1980--2018, demonstrating that active super El Ni\~no episodes (shaded intervals, 45 months = 9.66\% of record) account for the dominant portion of third-order covariance.}
\label{fig:jackknife}
\end{figure}

Table~\ref{tab:jackknife} and Figure~\ref{fig:jackknife} present the Jackknife experiment results. Excluding the 45 cumulative active months across the three super El Ni\~no episodes (1982--83: 16 months; 1997--98: 15 months; 2015--16: 14 months; representing exactly $9.66\%$ of the 466-month record, leaving $N = 421$ irregular baseline months) produces a pronounced collapse:
\begin{itemize}
    \item Univariate skewness $T_{EEE}$ drops by \textbf{73.5\%} under re-standardization (to $0.4159$) and reverses sign under global standardization ($-0.1800$).
    \item Mixed moment $T_{EEH}$ drops by \textbf{89.5\%} under re-standardization (to $0.1063$) and to $-0.1038$ under global standardization.
    \item Global tensor Frobenius norm drops by \textbf{66.1\%} under global standardization (from $3.2061$ to $1.0873$) and by \textbf{58.3\%} under re-standardization ($1.3356$).
\end{itemize}
Notably, evaluating the events individually demonstrates that the 1997--98 event exerts the largest single leverage ($-50.8\%$ norm change), followed by 1982--83 ($-32.8\%$), whereas excluding 2015--16 reduces the norm by only $-9.6\%$. This reflects the distinct longitudinal spatial structure of 2015--16, where warming was centered further westward toward the dateline (high Central Pacific index $C > +1.5\sigma$, but more moderate eastern index $E$), in contrast to the intense eastern Pacific cold tongue warming characterizing 1982--83 and 1997--98.

These results demonstrate that empirical third-order moments in the tropical Pacific are overwhelmingly concentrated in rare extreme episodes. This concentration is physically compatible with the atmospheric deep convective threshold ($\sim 28^\circ$C) in the cold tongue \citep{takahashi2016strong, takahashi2019theoretical, capotondi2020diversity} and oceanic non-linear dynamical heating \citep{an2004nonlinearity}. However, as noted by \citet{takahashi2019theoretical}, empirical subsampling across only three historical super El Ni\~no events is insufficient to unequivocally demonstrate an abrupt regime bifurcation over continuous higher-order non-linearities. The Jackknife thus serves as quantitative evidence of moment concentration in extreme episodes rather than a definitive proof of a discontinuous threshold.

Importantly, as noted in Section~2.2, removing these 45 months leaves three multi-month gaps in the observational series. Therefore, while this Jackknife subsample provides a rigorous test of static moment concentration, the irregular remaining sample ($N = 421$) cannot be treated as a continuous time series for lag-dependent dynamic analyses.

\subsection{Physical Event Composites and Oceanographic Distinctions}

\begin{figure}[!htbp]
\centering
\includegraphics[width=\textwidth]{fig07_composiciones_eventos_ENSO.pdf}
\caption{Event-aligned physical composites and phase-space trajectory, sharing common y-axis scaling across (a--c). (a) Super El Ni\~no composite (1982--83, 1997--98, 2015--16; $N=3$) showing standardized anomalies of $E$ (red), $C$ (orange), and eastern thermocline depth $D20$ ($H$, blue; shading denotes $\pm 1\sigma$), with individual event trajectories plotted as thin lines. (b) Moderate El Ni\~no composite (1987, 1991, 2002, 2009; $N=4$). (c) La Ni\~na composite for $N=4$ independent episodes (1988, 1998, 2007, 2010), aligned at the Central Pacific cold peak in $C$ (with 1998--2000 treated as a single prolonged multi-year event). (d) Phase-space trajectory in the $(E, D20)$ plane for the Super El Ni\~no composite from lag $-6$ to $+6$ months, with directional arrows tracing precursor local thermocline deepening ($t = -3$ mo), rapid eastern surface peak ($t = 0$), and subsequent rapid thermocline shoaling ($t = +3$ to $+6$ mo). Local eastern $D20$ measures boundary thermocline displacement, distinguished from basin-wide equatorial Warm Water Volume (WWV; \citealt{jin1997paradigm, meinen2000observations}).}
\label{fig:composites}
\end{figure}

Figure~\ref{fig:composites} shows that the third-order moment $\mathbb{E}[E(t) C(t-3) H(t-2)]$ is physically consistent with actual oceanic sequences rather than statistical artifacts:
\begin{itemize}
    \item \textit{Super El Ni\~no Composite (1982/83, 1997/98, 2015/16):} During the precursor phase, anomalous upper-ocean heat content accumulates rapidly: at $t = -3$ months, $C = +1.11\sigma$ and $H = +1.74\sigma$; at $t = -2$ months, $C = +1.29\sigma$ and $H = +2.54\sigma$, reaching $E = +3.09\sigma$ and $H = +3.38\sigma$ at peak eastern warming ($t = 0$). Following the peak, thermocline shoaling and cold tongue cooling proceed rapidly: by $t = +3$ months, the thermocline anomaly collapses to $H = +0.98\sigma$ (with $C = +0.82\sigma$), and by $t = +6$ months the eastern thermocline is exhausted ($H = -0.11\sigma$) while $C$ drops sharply into negative territory ($-0.48\sigma$).
    \item \textit{Moderate El Ni\~no Composite:} Eastern warming remains modest ($E = +0.48\sigma$ at $t = 0$), while Central Pacific warming is pronounced ($C = +1.43\sigma$ at $t = 0$ and $+1.46\sigma$ at $t = +1$ mo) and thermocline depth remains moderately elevated ($H = +1.74\sigma$ at $t=0$).
    \item \textit{La Ni\~na Composite:} Characterized by persistent shallow thermocline conditions ($H = -1.72\sigma$ at $t = 0$), cold Central Pacific anomalies ($C = -2.18\sigma$), and negative eastern SST ($E = -0.29\sigma$).
\end{itemize}

\textbf{Conceptual Oceanographic Nuance:} It is critical to distinguish local eastern $D20$ from the basin-wide Warm Water Volume (WWV) defined in the recharge oscillator paradigm \citep{jin1997paradigm, meinen2000observations, vialard2025recharge}. The classic recharge oscillator operates through the zonally averaged thermocline depth / upper-ocean heat content across the entire equatorial belt ($120^\circ\text{E}$--$80^\circ\text{W}$, $5^\circ\text{S}$--$5^\circ\text{N}$). Eastern equatorial $D20$ ($110^\circ\text{W}$--$95^\circ\text{W}$), by contrast, represents a local wave receiver that captures both the basin-wide recharge mode and the dynamical thermocline tilt mode driven by zonal wind stress anomalies ($\tau_x$). Hence, while precursor deepening of $D20$ at lag $-2$ months and its rapid collapse at lag $+3$ to $+6$ months reflect localized wave convergence and thermocline shoaling along the South American cold tongue, they should be understood as a surface-thermocline asymmetry diagnostic rather than direct equatorial basin-wide recharge. Furthermore, while the linear regression between $D20$, $E$, and $C$ accounts for $R^2 = 66.54\%$, this contemporaneous statistical fit reflects strong wave coupling but does not by itself prove pure linear wave kinematics nor preclude unresolved subsurface non-linear dynamics.

\section{Discussion}

\subsection{Physical Attribution: Atmospheric Thresholds, Oceanic NDH, and Subsurface Wave Coupling}
Our findings clarify the physical nature of ENSO non-linearity and position our results relative to recent dynamical frameworks:
\begin{enumerate}
    \item \textbf{Subsurface Wave Integration vs. Subsurface Non-Linearities:} The linear regression of eastern thermocline depth on SST indices yields $R^2 = 66.54\%$, reflecting the prominent role of linear equatorial wave dynamics in transferring wind-forced thermocline anomalies across the Pacific basin. However, we emphasize that a high contemporaneous linear regression does not by itself prove pure linear wave kinematics, nor does it preclude unresolved non-linear subsurface advection. Rather, it indicates that static third-order tensor projections are dominated by first-order thermocline tilt and wave convergence.
    \item \textbf{Dual Origin of Non-Linearity and Extreme Event Concentration:} The extreme non-linear behavior of $E$ arises from two established physical mechanisms: (i) the atmospheric deep convective threshold ($\sim 28^\circ$C) in the equatorial cold tongue \citep{takahashi2016strong, takahashi2019theoretical, capotondi2020diversity}, and (ii) oceanic non-linear dynamical heating (NDH; \citealt{an2004nonlinearity}), which amplifies warm extremes via vertical and horizontal temperature advection. While our Jackknife analysis shows that third-order statistics collapse dramatically (by $66.1\%$) when super El Ni\~no events are excluded, we concur with \citet{takahashi2019theoretical} that having only three historical extreme events limits empirical power to prove a discontinuous regime threshold over continuous non-linearities.
    \item \textbf{Contrast with Bispectral and Nonlinear Recharge Models:} In frequency space, \citet{pires2021bispectral} utilized bispectral analysis to identify quadratic phase coupling in ENSO indices, demonstrating that stochastic quadratic models capture substantial skewness. In physical space, \citet{han2026realistic} established that realistic ENSO simulation requires state-dependent non-linear damping in recharge oscillator dynamics. Our bi-lagged multilinear metric $\Delta T_{ECH}(+3, +2)$ complements these approaches by providing a compact time-domain diagnostic of the surface-thermocline asymmetry, distinguishing localized eastern thermocline response from basin-wide heat recharge.
    \item \textbf{Mathematical Implication:} Static polynomial tensor expansions are ill-suited to model this physics because extreme ENSO excursions operate as episodic, localized amplifications rather than continuous ambient multilinear turbulence.
\end{enumerate}

\subsection{Diagnostic Value of Bi-Lagged Moments for Climate Models}
While static third-order tensors provide little incremental information, the bi-lagged surface $T_{ECH}(\tau_1, \tau_2)$ and its time-reversal asymmetry $\Delta T_{ECH}$ offer powerful diagnostic benchmarks for evaluating coupled general circulation models (AOGCMs) in CMIP6 and CMIP7. As highlighted in recent comprehensive syntheses of the recharge oscillator \citep{vialard2025recharge, han2026realistic}, capturing realistic duration, discharge, and transition asymmetries between El Ni\~no and La Ni\~na remains a fundamental challenge for global climate models. Many state-of-the-art models suffer from persistent cold tongue biases, over-active central Pacific variability, and overly symmetric El Ni\~no/La Ni\~na transitions. The bi-lagged multilinear metric packages the phase progression of thermocline recharge ($\tau_2 = +2$ mo), precursor Central Pacific warming ($\tau_1 = +3$ mo), peak eastern warming ($t=0$), and post-peak transition into a compact scalar benchmark, providing an objective test for the temporal arrow of ENSO.

\section{Conclusions}

This study has presented a diagnostic evaluation of the trivariate $(E, C, D20)$ system over 1980--2018 ($N = 466$ months):
\begin{enumerate}
    \item \textbf{Static Multilinear Partitioning:} The static third-order moment tensor $\mathcal{T}^{(3)}$ is largely ($71.76\%$ squared Frobenius ratio) accounted for by linear thermocline wave coupling and univariate skewness ($T_{EEE} = 1.5708$, $T_{CCC} = -0.4430$). Geometric decomposition confirms that the linear reconstruction contributes $71.76\%$ ($\|\mathcal{T}_{\text{lin}}\|_F = 2.7160$), the non-orthogonal cross-term represents $21.82\%$ and the residual tensor norm represents $6.42\%$ ($\|\mathcal{T}_{\text{res}}\|_F = 0.8121$, with total norm $\|\mathcal{T}\|_F = 3.2061$). The general Gram-Schmidt identity confirms that $T_{EEC}$ is algebraically partitioned into linear correlation and Takahashi's quadratic curve ($m_{\text{res}} \equiv \gamma \text{Var}(q) = 0.1367$), without unprojected cubic components in the static phase plane.
    \item \textbf{Temporal Irreversibility and Arrow of Time:} The bi-lagged surface $T_{ECH}(\tau_1, \tau_2)$ identifies an isolated peak at $(\tau_1 = +3, \tau_2 = +2)$ months ($+0.4293$) with a time-reversal asymmetry of $\Delta T = +0.2405$. While an unconstrained grid search across 625 lags yields $p = 0.44$ due to random field extreme-value fluctuations, multivariate Fourier surrogates ($p = 0.0105$), trivariate Gaussian VAR(2) models ($p = 0.0295$), and a dynamic non-linear coupled recharge model ($p = 0.0140$) reject linear/quadratic benchmarks when constrained to the physical wave-transit timescale $(\tau_C = +3, \tau_H = +2)$. Retrospective temporal cross-validation on 2000--2018 ($N=228$) confirms this positive asymmetry ($\Delta T = +0.3227$, 95\% bootstrap CI: $[+0.0319, +0.5125]$, Fourier surrogate $p = 0.0395$, dynamic quadratic null model $p = 0.0065$ [one-tailed], $p=0.0150$ [two-tailed, $p < 0.05$]). Physical event composites confirm that this asymmetry reflects the localized surface-thermocline phase progression compatible with recharge-discharge dynamics, while distinguishing local eastern thermocline shoaling from basin-wide Warm Water Volume (WWV; \citealt{jin1997paradigm, vialard2025recharge}).
    \item \textbf{Moment Concentration and Sensitivity in Extreme Episodes:} Excluding the three super El Ni\~no events (45 active months; $9.66\%$ of the record) produces a $66.1\%$ reduction in the global tensor Frobenius norm under global standardization (from $3.2061$ to $1.0873$) and a $58.3\%$ reduction under re-standardization ($1.3356$). This demonstrates the high leverage and sensitivity of empirical higher-order moments to rare extreme episodes, consistent with atmospheric convective thresholds \citep{takahashi2016strong, takahashi2019theoretical} and oceanic non-linear dynamical heating \citep{an2004nonlinearity}, though historical sample sizes remain insufficient to definitively separate threshold bifurcations from continuous higher-order non-linearities.
    \item \textbf{Predictive Falsification:} Incorporating bilinear interaction terms fails to improve out-of-sample forecast skill and produces slightly higher RMSE across all lead horizons (+0.81\% at 1 mo, +1.61\% at 3 mo, +1.82\% at 6 mo), failing to show empirical predictive utility.
\end{enumerate}

Overall, trivariate moment tensor analysis provides a robust, statistically significant diagnostic of surface-thermocline temporal asymmetry under physically motivated timescales, offering an objective benchmark for climate model fidelity, while its physical interpretation reflects known wave and convective mechanisms rather than autonomous multilinear physics.

\section*{Acknowledgments and Data Availability}
This study was conducted within the research framework of the ARGO-CROSS Project at the Instituto Geof\'isico del Per\'u (IGP). Primary observational datasets are publicly accessible: NOAA ERSSTv5 (\url{https://www.ncei.noaa.gov/products/extended-reconstructed-sst}), TAO/TRITON moored buoy array data from NOAA PMEL (\url{https://www.pmel.noaa.gov/tao/}), and ECMWF ORAS5 ocean reanalysis data (\url{https://www.ecmwf.int/en/forecasts/dataset/ocean-reanalysis-system-5}). All Python processing scripts (\texttt{audit\_01} to \texttt{audit\_07}, \texttt{generate\_all\_final\_figures.py}), processed NetCDF files, synchronized CSV data tables (\texttt{tabla\_datos\_sincronizados\_1980\_2018.csv}), Monte Carlo simulation routines ($B=2000$, random seed 42), and LaTeX source files are archived and will be made publicly available upon formal acceptance at the dedicated GitHub repository (\url{https://github.com/mandrade-igp/argo-cross-enso-tensor}) and permanent Zenodo archive. During peer review, all reproducible scripts, data tables, and figures are fully accessible to reviewers and editors via the companion Overleaf repository mirror (Project ID: \texttt{6ae69499583b59198223570c}) and upon direct request to the corresponding author (\texttt{mandrade@igp.gob.pe}). Numerical analyses were executed under Python 3.11 with standard scientific packages (\texttt{numpy 1.26}, \texttt{scipy 1.12}, \texttt{statsmodels 0.14}, and \texttt{matplotlib 3.8}).

\bibliographystyle{plainnat}
\bibliography{references_final}

\end{document}
